\section{Counting the balanced port profiles}
\label{sec:counting}

The port model of Proposition~\ref{prop:port-model} and the evaluator
of Proposition~\ref{prop:weight-evaluator} give a sum of positive weights
on polynomially short ranges whose individual terms we can approximate.
We now approximate the entire sum.
The factorial encoding below is essential: its transversals have exactly
the mass of the unlabelled port profiles.  The two resulting binary weights
allow us to use the transport and trace estimates of the common-base
framework, in the two-weight form proved in~\cite{CB}.

Write $\mathcal P$ for the positive-length ports and $d_v$ for the length
of port $v\in\mathcal P$.  Thus $p=\sum_v d_v$, and the weight $h$ of
Proposition~\ref{prop:port-model} is positive on the entire slice
\[
 \mathcal B_p=\left\{(a,b):0\leq a_v,b_v\leq d_v,
                  \quad\sum_v(a_v+b_v)=p\right\}.
\]
All count coordinates in this section are integers.  The quantity to be
estimated is
\begin{equation}
 \mathcal Z_{\rm port}
   =\sum_{\substack{(a,b)\in\mathcal B_p\\a_v+b_v=d_v\ (v\in\mathcal P)}}
           h(a,b),
 \qquad
 \left|\log\frac{\mathcal Z_{\rm port}}{JZ}\right|
        \leq100N^{-7}.
 \label{eq:counting-target}
\end{equation}
Length-zero ports have already been fixed in $h$.  If $p=0$, the sum has
one term; we evaluate it directly at the end of the proof.  Until then
assume $p\geq1$.

\subsection{Two factorial weights and their signatures}

For every port $v$, introduce $d_v$ labelled pairs, each containing an
$a$-element and a $b$-element.  Let $\mathcal E$ be this ground set of
$2p$ elements and number its pairs $1,\ldots,p$.  For every $p$-subset
$S\subseteq\mathcal E$, its counts $(a(S),b(S))$ belong to
$\mathcal B_p$.  Given a positive count weight $k$ on this slice, define
\begin{equation}
 f(S)=k(a,b)\prod_{v\in\mathcal P}\frac{a_v!b_v!}{d_v!},
 \qquad
 \widetilde f(S)=k(a,b)
       \prod_{v\in\mathcal P}\frac{(d_v-a_v)!(d_v-b_v)!}{d_v!}.
 \label{eq:counting-lifts}
\end{equation}
Both weights are defined and positive on \emph{every} $p$-subset, including
subsets with several empty or full pairs.

Let $\mathcal A$ consist of the transversals, which select exactly one
element from each pair.  For distinct labels $i,l$, let $\mathcal A_{il}$
consist of the subsets with pair $i$ empty, pair $l$ full, and every
other pair singly occupied.  On a transversal the two factors in
\eqref{eq:counting-lifts} coincide and equal
$\prod_v\binom{d_v}{a_v}^{-1}$.  Exactly
$\prod_v\binom{d_v}{a_v}$ transversals realize any balanced profile.
Consequently
\begin{equation}
 C:=\sum_{S\in\mathcal A}f(S)
   =\sum_{S\in\mathcal A}\widetilde f(S)
   =\sum_{a+b=d}k(a,b).
 \label{eq:counting-multiplicity}
\end{equation}
This identity concerns transversals, rather than all subsets with the
same coordinate counts.

For a positive weight on all $r$-subsets of a finite set, the
\emph{addition property} means that, after any $r-2$ elements have
been fixed included, the matrix of weights obtained by adding two
distinct further elements, with zero diagonal, has at most one positive
eigenvalue.  It is vacuous for $r<2$.  The \emph{omission property} is
the addition property of the complement-set weight.

\begin{lemma}
\label{lem:counting-lifts}
If $k$ has both integer signatures of Definition~\ref{def:integer-signatures},
then $f$ has the omission property and $\widetilde f$ has the addition
property.  Both properties persist under fixing included or excluded
elements and summing disjoint labelled pairs at occupancy one, on each
resulting full valid-degree domain.  Moreover, with $q=(p+1)^2$,
\begin{equation}
 q^{-1}f(S)\leq\widetilde f(S)\leq qf(S)
 \quad\left(S\in\mathcal A\cup\bigcup_{i\ne l}\mathcal A_{il}\right).
 \label{eq:counting-comparison}
\end{equation}
\end{lemma}

\begin{proof}
We give the factorial cancellation in
\cite[Lemma~4.2, ``Signatures of the two lifts'']{CB}, since it also
explains the roles of the two weights.  In an omission test for $f$,
the display has $p+2$ elements.  Group these by count coordinate and let
$n_u$ be the number in group $u$.  On the subspace of vectors constant
within each group, the ordered-copy multiplicities $n_un_w$ for $u\ne w$
and $n_u(n_u-1)$ for $u=w$ cancel the factorial reductions from removing
two elements.  The quadratic form becomes
\[
 \frac{\prod_u n_u!}{\prod_v d_v!}
       \sum_{u,w}k(n-e_u-e_w)t_ut_w,
\]
where invalid removals contribute zero.  Its positive index is at most
one by the integer removal signature.  Vectors summing to zero within
each group have nonpositive quadratic form and are orthogonal for this
form to the group-constant subspace: within a group the matrix has
constant nonnegative off-diagonal entries and zero diagonal.  This
proves the omission assertion.  For $\widetilde f$, group the elements
available after fixing $p-2$ inclusions.  The remaining-capacity
factorials cancel the same multiplicities, leaving a positive multiple
of the integer addition matrix.  This proves the addition assertion.
The closure assertions are precisely
\cite[Lemma~4.1, ``Binary contractions'']{CB}; their hypothesis of
positivity on the full domain has been checked here.

For the comparison, if the empty and full pair are in the same port,
all ports remain balanced and the ratio is one.  Otherwise the port
containing the hole has $a+b=d-1$, and contributes $(a+1)(b+1)$
to $\widetilde f/f$.  The port containing the full pair has
$a+b=d+1$ and contributes $1/(ab)$.  Both positive integer products
lie between $1$ and $(p+1)^2$.  Their quotient proves
\eqref{eq:counting-comparison}.
\end{proof}

\subsection{The transport input and the chain it controls}

We state the precise consequence of the companion transport argument
that will be used.  This isolates the long transport proof from the
algorithm and makes its hypotheses available for verification.
Fix $f,\widetilde f$ as above and put
$c_{il}=\sum_{S\in\mathcal A_{il}}f(S)$.  These numbers are positive.
For positive multipliers $w_{il}$ define, on
$\mathcal X=\mathcal A\cup\bigcup_{i\ne l}\mathcal A_{il}$,
\[
 \lambda(S)=
 \begin{cases}f(S),&S\in\mathcal A,\\
               w_{il}f(S),&S\in\mathcal A_{il},\end{cases}
 \quad
 \Lambda=\sum_{S\in\mathcal X}\lambda(S),
 \quad \pi(S)=\lambda(S)/\Lambda,
 \quad \mu(S)=f(S)/C\quad(S\in\mathcal A).
\]
Call the multipliers \emph{good} when
\begin{equation}
 \frac{C}{4c_{il}}\leq w_{il}\leq\frac{4C}{c_{il}}
 \qquad(i\ne l).
 \label{eq:counting-good}
\end{equation}
In this case
\begin{equation}
 \Lambda\leq5p^2C,\qquad
 \pi(\mathcal A)\geq(5p^2)^{-1},\qquad
 \pi(\mathcal A_{il})\geq(20p^2)^{-1}.
 \label{eq:counting-type-mass}
\end{equation}

The chain holds with probability $1/2$.  Otherwise it chooses independently
a uniform element of $S$ and a uniform element of $\mathcal E\setminus S$,
and proposes their exchange.  A proposal outside $\mathcal X$ is rejected;
a proposal $S'\in\mathcal X$ is accepted with probability
$\min(1,\lambda(S')/\lambda(S))$.  Write $P$ for this kernel.  A defect
state can reach a transversal by moving an element of its full pair to
its empty pair, and transversals are connected by flips within pairs.
Thus positivity makes $P$ irreducible.  It is lazy and reversible, with
Dirichlet form
\begin{equation}
 \mathcal E_P(g)=\langle g,(I-P)g\rangle_\pi
   =\frac{\mathfrak D(g)}{2p^2\Lambda},\qquad
 \mathfrak D(g)=\sum_{\{S,S'\}}
      \min(\lambda(S),\lambda(S'))(g(S)-g(S'))^2,
 \label{eq:counting-energy}
\end{equation}
where the sum is over unordered exchange edges.
Let $\mathcal Qg$ agree with $g$ at each transversal and replace its
values on each defect class by their conditional $\pi$-mean.

\begin{proposition}[Companion transport and simulation input]
\label{prop:counting-transport}
Suppose $f$ is positive on all $p$-subsets and has the omission property,
$\widetilde f$ is positive on that domain and has the addition property,
the weights agree on $\mathcal A$, and
\eqref{eq:counting-comparison} holds with $q\geq1$.
For good multipliers put
\begin{equation}
 C_{\rm p}=p(1+8pq^2),\qquad
 C_{\rm d}=8+32pq^3,\qquad
 K_{\rm obs}=10p^4(9C_{\rm p}+2C_{\rm d}).
 \label{eq:counting-constants}
\end{equation}
Then for every real function $g$ on $\mathcal X$,
\begin{equation}
 C\operatorname{Var}_\mu(g)\leq C_{\rm p}\mathfrak D(g),\qquad
 \operatorname{Var}_\pi(\mathcal Qg)
 \leq(9C_{\rm p}+2C_{\rm d})\frac{\mathfrak D(g)}C
 \leq K_{\rm obs}\mathcal E_P(g).
 \label{eq:counting-variance}
\end{equation}
If $\mathcal QG=G$ and an initial law satisfies $\nu\leq D\pi$
pointwise, then $m\geq1$ consecutive observations satisfy
\begin{equation}
 \mathbb E_\nu\left[
  \left(m^{-1}\sum_{t=0}^{m-1}G(S_t)-\mathbb E_\pi G\right)^2\right]
 \leq\frac{2D K_{\rm obs}}m\operatorname{Var}_\pi(G).
 \label{eq:counting-time-average}
\end{equation}
The trace on $\mathcal A$, counting every strictly positive return as
one step, is lazy and reversible with stationary law $\mu$ and inverse
spectral gap at most $2p^2C_{\rm p}\leq K_{\rm obs}$.  In particular,
\begin{equation}
 \left|\frac{P_{\rm tr}^t(x,y)}{\mu(y)}-1\right|
 \leq\mu_{\min}^{-1}\exp(-t/K_{\rm obs}).
 \label{eq:counting-trace-mixing}
\end{equation}
If a trace starts from $\nu\leq D\mu$, the expected cost of $t$ trace
steps in underlying transitions is at most $tD/\pi(\mathcal A)$;
at each deterministic trace index its law is still at most $D\mu$.
\end{proposition}

\begin{proof}[Source and applicability]
These are \cite[Propositions~5.2 and~5.4 and Lemmas~5.5--5.6]{CB},
with $q$ left as the comparison parameter in their proofs.  Their
structural inputs are Lemma~5.1 (one-hole transport) and Lemma~4.4
(pair inequalities for defect masses).  We spell out why the stated
hypotheses are the complete interface.  For a partial transversal
assignment $\sigma$, let $z_\sigma$ be its transversal mass and let
$\widetilde c_{il}^{\sigma}$ be the mass of its extensions in
$\mathcal A_{il}$ under $\widetilde f$.  Binary contraction, including
the full-index cubic contraction clause of \cite[Lemma~4.1]{CB},
and the addition property give
\[
 4\widetilde c_{il}^{\sigma}\widetilde c_{li}^{\sigma}\leq z_\sigma^2,
 \qquad
 \widetilde c_{il}^{\sigma}\widetilde c_{lk}^{\sigma}
       \leq z_\sigma\widetilde c_{ik}^{\sigma}
 \quad(i,l,k\text{ distinct and unassigned}).
\]
Agreement on transversals identifies $z_\sigma$ for the two weights.
The omission property supplies one-hole transport for $f$, including
its auxiliary displays outside $\mathcal X$; this is why positivity
and the signature were required on all $p$-subsets.
Comparison on actual defect classes costs $q^2$ in the transversal
variance bound and $q^3$ in the defect-mean bound, giving exactly
$C_{\rm p},C_{\rm d}$ above.  The remaining observable and trace
arguments use only these inequalities, reversibility, and laziness.
They impose no condition on the origin of the count weight.
This is the two-lift version of the transport, observable-Poincare,
and trace-restart framework of~\cite[Sections~4--6]{Parent}.
\end{proof}

\subsection{Power annealing and exact initialization}

Lemma~\ref{lem:weight-bounds} gives $|\log h|\leq N^{60000}$, and
Proposition~\ref{prop:port-model} gives both integer signatures for
$h^s$ at every $0\leq s\leq1$.
Choose
\begin{equation}
 U=N^{61000},\qquad H_o=10U,\qquad
 k_j=(2^{-U}h)^{j/H_o}\quad(0\leq j\leq H_o),
 \qquad K_o=10(p+U+1).
 \label{eq:counting-schedule}
\end{equation}
Apply \eqref{eq:counting-lifts} to $k_j$, obtaining $f_j,\widetilde f_j$,
and denote their transversal and defect masses by $C_j,c_{il}(j)$.
Lemma~\ref{lem:counting-lifts} verifies every structural hypothesis of
Proposition~\ref{prop:counting-transport} at every phase, with
$q=(p+1)^2$.  In particular
\begin{equation}
 C_{H_o}=2^{-U}\mathcal Z_{\rm port}.
 \label{eq:counting-final-scale}
\end{equation}

\begin{lemma}[Initialization and adjacent phases]
\label{lem:counting-initialization}
For every $p$-subset and $0\leq j<H_o$,
\begin{equation}
 \frac12\leq\frac{f_{j+1}(S)}{f_j(S)}\leq1.
 \label{eq:counting-adjacent}
\end{equation}
Thus $C_j/c_{il}(j)$ changes by a factor in $[1/2,2]$ between
consecutive phases, and $\mu_j=f_j/C_j$ on $\mathcal A$ satisfies
$\mu_j\leq2\mu_{j+1}$.  Moreover
\begin{equation}
 \min_{S\in\mathcal A}\mu_j(S)
       \geq2^{-(H_o+p+K_o)}.
 \label{eq:counting-minimum}
\end{equation}
The initial masses are exact rational numbers computable in polynomial
bit time, and $\mu_0$ can be sampled by independent uniform port counts
followed by uniform subsets of the pairs realizing those counts.
\end{lemma}

\begin{proof}
The logarithm of $2^{-U}h$ lies between
$-U\log2-N^{60000}$ and $-U\log2+N^{60000}<0$.
Dividing by $H_o=10U$ proves \eqref{eq:counting-adjacent}.
Summing that bound within each type proves the multiplier comparison
and domination.  On a transversal the factorial factor is at least
$2^{-p}$, while $k_j\geq2^{-H_o}$.  Also
\[
 C_j\leq C_0=\prod_{v\in\mathcal P}(d_v+1)\leq2^p\leq2^{K_o},
\]
which proves \eqref{eq:counting-minimum}.

Here are explicit formulas for the initializer, as in
\cite[Lemma~7.1, ``Exact initialization'']{CB}.  Fix $i\ne l$.
For each port put
\[
 u_v=\mathbf1_{\{i\text{ lies in }v\}},\quad
 v_v=\mathbf1_{\{l\text{ lies in }v\}},\quad
 m_v=d_v-u_v-v_v.
\]
There are $m_v$ singly occupied pairs left in this port, so
\begin{equation}
 c_{il}(0)=\prod_{v\in\mathcal P}
 \left[\frac1{d_v!}\sum_{k=0}^{m_v}
      \binom{m_v}{k}(k+v_v)!(m_v-k+v_v)!\right].
 \label{eq:counting-initial-defect}
\end{equation}
All factorial arguments are at most $p$.  If $i,l$ are in the same
port, that port has at least two pairs, so $m_v\geq0$ also in this case.
Each bracket equals, respectively,
$d_v+1$, $1$, $(d_v+1)(d_v+2)/6$, or $(d_v+1)/6$ according as
the port contains neither designated pair, only $i$, only $l$, or both.
Writing $v(i)$ for the port containing pair $i$, we can therefore use
the particularly simple exact initial multipliers
\begin{equation}
 w_{il}^{(0)}=\frac{C_0}{c_{il}(0)}=
 \begin{cases}
 6,&v(i)=v(l),\\[1mm]
 \displaystyle\frac{6(d_{v(i)}+1)}{d_{v(l)}+2},&v(i)\ne v(l).
 \end{cases}
 \label{eq:counting-initial-multipliers}
\end{equation}
The finite sums or these simplified expressions give polynomial bit
bounds, including all $p(p-1)$ initial multipliers.
Finally choose $a_v$ uniformly from $\{0,\ldots,d_v\}$ and choose
a uniform $a_v$-subset of that port's pairs to use their $a$-elements.
The probability of a particular transversal is
$\prod_v[(d_v+1)\binom{d_v}{a_v}]^{-1}=f_0(S)/C_0$.
\end{proof}

\subsection{Capped restarts and ratio estimates}

The next experiment uses exact real transition probabilities only for
its analysis.  Its finite implementation follows in the next subsection.
Use \eqref{eq:counting-constants} with $q=(p+1)^2$ and choose
\begin{equation}
 \begin{aligned}
 \varepsilon_0&=\varepsilon/100,
 &\eta&=\frac{\varepsilon_0}{200(H_o+1)},\\
 \tau&=2K_{\rm obs}(H_o+p+K_o+5),
 &B_{\rm ret}&=2000p^2(H_o+1)^2\tau,\\
 N_{\rm av}&=\left\lceil
       \frac{10^8p^8K_{\rm obs}(H_o+1)}{\eta^2}\right\rceil.
 \end{aligned}
 \label{eq:counting-caps}
\end{equation}
These are integers except for the tolerances.  We write $P_{j,w}$ for
the chain at phase $j$ and multiplier collection $w$.
Starting with \eqref{eq:counting-initial-multipliers}, perform the
following operations for $j=0,\ldots,H_o-1$.

\begin{enumerate}
\item Draw a fresh transversal from $\mu_0$.  For $a=1,\ldots,j$,
run $\tau$ trace steps of $P_{a,w^{(a)}}$, continuing from the
transversal reached at the preceding stage.  Each trace step ends at
the next transversal after at least one underlying transition;
holding at a transversal counts as a return.  Allow at most
$B_{\rm ret}$ underlying transitions across the entire restart.
If another is needed, abort and return zero.  When $j=0$ the fresh
$\mu_0$ draw itself is the restart endpoint.

\item From that endpoint take $N_{\rm av}$ consecutive observations
of $P_{j,w^{(j)}}$, including the initial state.  Label the transversal
type by $0$ and the other types by $il$.  Let $n_T$ be the number of
observations of type $T$, and set $\widehat p_T=n_T/N_{\rm av}$.
Also compute the empirical mean $\widehat u_j$ of
\[
 G_j(S)=\mathbf1_{\mathcal A}(S)\frac{f_{j+1}(S)}{f_j(S)}.
\]
If any type count is zero, abort and return zero.

\item Set $\widehat R_j=\widehat u_j/\widehat p_0$ and, if
$j+1<H_o$, store the exact rational updates
\begin{equation}
 w_{il}^{(j+1)}=w_{il}^{(j)}\frac{n_0}{n_{il}}.
 \label{eq:counting-update}
\end{equation}
\end{enumerate}
All randomness in a phase, including its initial sample, is fresh.
If no abort occurs, return
$Y_{\rm ref}=C_0\prod_{j=0}^{H_o-1}\widehat R_j$.
The multipliers in a restart have already been stored by earlier phases.

\begin{proposition}
\label{prop:counting-reference}
With probability at least $1-1/50$, this experiment does not abort,
every stored multiplier collection is good at its own phase, and
\begin{equation}
 \left|\log\frac{Y_{\rm ref}}{C_{H_o}}\right|\leq4H_o\eta.
 \label{eq:counting-reference-error}
\end{equation}
\end{proposition}

\begin{proof}
This is the reference experiment of \cite[Section~7.2 and
Proposition~7.2]{CB}.  We include its probabilistic argument to specify
the treatment of the restart caps.  Fix a history of successful phases
before $j$.  Its stored multipliers are now fixed and good.  For this
one phase, consider a virtual restart that ignores the transition cap,
followed by its observation trajectory.  At any trace stage $a$,
\eqref{eq:counting-minimum} and \eqref{eq:counting-trace-mixing} bound
the pointwise relative error after $\tau$ steps by
\[
 2^{H_o+p+K_o}\exp(-2(H_o+p+K_o+5))<1.
\]
The endpoint law is thus at most $2\mu_a$.  Stage one starts at
$\mu_0\leq2\mu_1$, and each later stage starts at a law bounded
by $2\mu_{a-1}\leq4\mu_a$.  Trace stationarity preserves the bound
$4\mu_a$ at every deterministic trace index in that stage.
By the return-cost assertion and \eqref{eq:counting-type-mass},
each such trace step has expected cost at most $20p^2$.
There are at most $H_o\tau$ trace steps in a restart.  Markov's
inequality therefore bounds its cap-failure probability by
\begin{equation}
 \frac{20p^2H_o\tau}{B_{\rm ret}}
       \leq\frac1{100(H_o+1)}.
 \label{eq:counting-cap-error}
\end{equation}

The virtual endpoint has law at most $2\mu_j$, hence at most
$10p^2\pi_{j,w^{(j)}}$ on $\mathcal X$.  Under this stationary law,
write $p_T$ for the type probabilities and $u_j=\mathbb E_\pi G_j$.
Direct summation gives
\begin{equation}
 \frac{u_j}{p_0}=\frac{C_{j+1}}{C_j},\qquad
 w_{il}^{(j)}\frac{p_0}{p_{il}}=\frac{C_j}{c_{il}(j)}.
 \label{eq:counting-observables}
\end{equation}
The type indicators and $G_j$ are fixed by $\mathcal Q$, take values
in $[0,1]$, and have means at least $(20p^2)^{-1}$, by
\eqref{eq:counting-type-mass} and \eqref{eq:counting-adjacent}.
The mean squared error of each average is at most
$20p^2K_{\rm obs}/N_{\rm av}$ by
\eqref{eq:counting-time-average}.  There are
$p(p-1)+2\leq2p^2$ observables.  Markov's inequality applied to
their squared errors and a union bound show that the probability
some empirical mean has relative error greater than $\eta$ is at most
\begin{equation}
 \frac{16000p^8K_{\rm obs}}{N_{\rm av}\eta^2}
       \leq\frac1{100(H_o+1)}.
 \label{eq:counting-observation-error}
\end{equation}
These estimates concern the uncapped virtual phase.  On the event
that its cost meets the cap it agrees with the capped phase.
Failure of the latter is contained in the union of the two failure
events just bounded; no endpoint law is conditioned on meeting a cap.

On empirical success every type count is positive.  Equation
\eqref{eq:counting-observables} shows that the update estimates the
current ideal multiplier directly, with factor error at most
$(1+\eta)/(1-\eta)\leq2$.  Its ideal value changes by at most a
further factor two for the next phase, so
\eqref{eq:counting-good} holds there.  The same empirical bounds give
\[
 \left|\log\frac{\widehat R_j}{C_{j+1}/C_j}\right|\leq4\eta.
\]
The initial multipliers are ideal.  Summing the probabilities of the
first unsuccessful phase bounds total failure by
$2H_o/[100(H_o+1)]<1/50$.  On success the ratios telescope and
their log errors sum to \eqref{eq:counting-reference-error}.
\end{proof}

\subsection{Implementation with bounded fair-bit computation}

To turn the experiment into an algorithm, put
\begin{equation}
 J_*=100(H_o+1)(B_{\rm ret}+N_{\rm av}+p+2n+1),
 \qquad \xi=\frac{\eta}{10000(J_*+1)}.
 \label{eq:counting-bit-parameters}
\end{equation}
To evaluate $k_j$, call Proposition~\ref{prop:weight-evaluator} freshly
for $h^{j/H_o}$ with relative tolerance $\xi/4$ and failure probability
$\xi$.  Approximate the known factor $2^{-Uj/H_o}=2^{-j/10}$
deterministically to relative error $\xi/4$ by rational bisection of
the tenth root of $2^{-j}$.  Its endpoints have $O(j+\log\xi^{-1})$
bits, so this is bounded polynomial arithmetic.  The product is a
positive rational evaluation of $k_j$ with relative error at most
$\xi$ and failure probability at most $\xi$.  Multiply by the exact
factorial factor to evaluate $f_j$.  Use fresh
evaluations for the two weights in each valid proposed Metropolis
ratio and in each positive observation of $G_j$.  The multipliers
and type-count updates remain exact rationals.  Cap each acceptance
ratio at one and implement it with a downward-rounded dyadic threshold
using $\lceil\log_2(1/\xi)\rceil$ bits.  The half-hold uses one exact
fair bit.

Here is a bounded implementation of every other finite uniform draw.
To draw from $\{0,\ldots,v-1\}$, choose a uniform $b$-bit integer
$Y$ with $2^b\geq v/\xi$ and return $\lfloor vY/2^b\rfloor$.
Each output has either $\lfloor2^b/v\rfloor$ or
$\lceil2^b/v\rceil$ preimages, so its total variation distance from
uniform is at most $v/2^b\leq\xi$.
For a subset of size $a_v$ use
$v=\binom{d_v}{a_v}$ and decode its index by successive binomial
inclusion and exclusion counts.  This takes at most $d_v$ decisions.
Thus initialization needs two finite draws per port, and each
proposal needs two finite draws.  No rejection loop is used.

\begin{lemma}
\label{lem:counting-bits}
The implemented experiment has deterministic polynomial bounds on
oracle calls, bit operations, and output length, on every execution.
It can be coupled to the capped reference experiment so that, except
with probability at most $20J_*\xi$, all states, type counts,
multiplier updates, and cap decisions coincide, every weight evaluation
meets its tolerance, and each positive observed ratio differs from
its reference value by log error at most $4\xi$.
\end{lemma}

\begin{proof}
There are at most $T_{\max}=H_o(B_{\rm ret}+N_{\rm av})$
underlying transitions.  Each requires at most two weight calls and
four outer draws (hold, two proposal choices, and acceptance).
The observations require at most $2H_oN_{\rm av}$ further calls,
and the restarts at most $2H_o|\mathcal P|$ initialization draws.
Including transition counts, the number of these elementary actions
is at most
\[
 7T_{\max}+2H_oN_{\rm av}+2H_o|\mathcal P|<J_*.
\]
Counter initialization and arithmetic on all defect multipliers add
only deterministic polynomial work.

For the coupling, expose the joint past immediately before each
call or draw.  Until a discrepancy, the states and multipliers agree.
The new evaluation input is then fixed by that past, and fresh bits
give conditional evaluation failure probability at most $\xi$.
This uses no conditioning on future successful evaluations.
Conditional on accurate evaluations, a computed weight ratio $r'$
and its exact value $r$ satisfy
\[
 \frac{1-\xi}{1+\xi}\leq\frac{r'}r
       \leq\frac{1+\xi}{1-\xi},\qquad
 |\min(1,r')-\min(1,r)|\leq\frac{2\xi}{1-\xi}<4\xi.
\]
Dyadic rounding contributes at most another $\xi$.
Couple reference and implemented acceptance decisions with a common
uniform variable.  Formally, the algorithm's uniform $b$-bit integer
can be supplemented, only in the proof, by an independent uniform
point inside its dyadic interval.  The resulting continuous uniform
variable gives the exact reference marginal and realizes the stated
probability of disagreement.  Couple each other finite draw at its
total variation cost, at most $\xi$.  A union bound over the action
budget, including call failures, gives the generous bound $20J_*\xi$.
Matching paths have identical integer type counts; consequently
\eqref{eq:counting-update} gives exactly identical stored multipliers,
even though the ratio observations differ.  Trace returns and cap
decisions also agree.  Successful positive ratios have log error at
most $4\xi$, as claimed.  This is the bounded coupling argument of
\cite[Lemma~7.3 and Section~7.3]{CB}, with the evaluator furnished here.

For the unconditional resource bound, $|\mathcal P|\leq2n$ and
$d_v\leq N^{21}W=N^{20021}$, so $p\leq2nN^{20021}$.
All numerical quantities $U,H_o,K_o,q,C_{\rm p},C_{\rm d},K_{\rm obs}$,
the displayed caps, and $\xi^{-1}$ are bounded by fixed polynomials
in $N$ and $\varepsilon^{-1}$.  Their exponents are absolute constants.
Every state has a $2p$-bit representation; its count profile and type
are recovered by a bounded scan.  Factorials and binomial coefficients
have $O(p\log(p+1))$ bits and are computed by bounded integer arithmetic.
Each weight evaluation has deterministic polynomial work and output
length by Proposition~\ref{prop:weight-evaluator}, even on its failure
event.  The scale $2^{-Uj/H_o}$ and rational powers have polynomial
descriptions and evaluation precision; restoring $2^U$ ultimately
requires only $U+1$ bits.

The sizes of the stored multipliers are controlled without any accuracy
assumption.  On a nonaborting history,
\[
 w_{il}^{(a)}=w_{il}^{(0)}
       \prod_{j=0}^{a-1}\frac{n_0^{(j)}}{n_{il}^{(j)}},
 \qquad1\leq n_T^{(j)}\leq N_{\rm av}.
\]
Thus each numerator and denominator has at most its initial bit length
plus $H_o\lceil\log_2(N_{\rm av}+1)\rceil$ bits, up to a fixed
additive constant, even without fraction reduction.  A zero count
causes an abort before division.  Acceptance thresholds therefore
have polynomial-length operands on every history.  Adding at most
$N_{\rm av}$ positive rational observations using a product denominator,
and multiplying at most $H_o$ phase estimates, also creates only
polynomial-length integers.  These estimates include the cost of
writing the rational output.  All outer loops have the specified caps,
all evaluator loops are bounded, and every random draw uses finitely
many fair bits.  The resource conclusions are independent of successful
coupling, good multipliers, and statistical accuracy.
\end{proof}

\subsection{Completion of the approximation}

\begin{proof}[Proof of Theorem~\ref{thm:main}]
First perform the exact feasibility test from Section~\ref{sec:polyhedra}.
Its integrality guarantee makes an empty real intersection equivalent
to $Z=0$, in which case the algorithm returns zero on every execution.
If $R=0$, the only feasible vector is zero and the algorithm returns
one; this includes the promised case $n=0$.  In all other feasible
cases perform the partition and port constructions proved above.

If $p\geq1$, let $Y$ be the implemented ratio product.  By
Proposition~\ref{prop:counting-reference} and
Lemma~\ref{lem:counting-bits}, outside an event of probability at most
$1/50+20J_*\xi<1/4$, the run does not abort and
\[
 \left|\log\frac{Y}{C_{H_o}}\right|
       \leq4H_o(\eta+\xi)<\varepsilon/100.
\]
Indeed the multiplicative bounds on each positive observation persist
under summation, while the type counts agree exactly.  Restore the
scale by setting $\widehat{\mathcal Z}_{\rm port}=2^UY$.
If $p=0$, evaluate its sole count profile directly with relative
tolerance $\varepsilon/100$ and failure probability $1/16$; then
$|\log(\widehat{\mathcal Z}_{\rm port}/\mathcal Z_{\rm port})|
\leq\varepsilon/50$ on success.  This branch invokes none of the
$p$-dependent chain parameters.  When $p=1$, the chain proof already
applies: there are no defect classes or multipliers and every state is
a transversal.

Obtain the deterministic positive rational estimate $\widehat J$ from
Proposition~\ref{prop:weight-evaluator}, with
$|\log(\widehat J/J)|\leq N^{-7}$, and output, for this run,
$\widehat Z=\widehat{\mathcal Z}_{\rm port}/\widehat J$.
Aborted runs output zero.  On the success event in either branch,
\eqref{eq:counting-target} gives
\[
 \left|\log\frac{\widehat Z}{Z}\right|
       \leq\varepsilon/50+101N^{-7}<\varepsilon/10.
\]
Here $N\geq10^{10}$ and $N\geq\varepsilon^{-1}$, so
$101N^{-7}\leq101\varepsilon N^{-6}<\varepsilon/100$.
For $0<\varepsilon<1$, $e^{-\varepsilon/10}\geq1-\varepsilon$
and $e^{\varepsilon/10}\leq1+\varepsilon$; hence a single run has
the required relative accuracy with probability greater than $3/4$.

For the requested failure probability, let
$s_\delta=\lceil\log_2(\delta^{-1})\rceil$ and take
$m_\delta=10s_\delta+1$ independent runs, returning their median.
If $F$ is their number of failures, independence implies
$\mathbb E2^F\leq(5/4)^{m_\delta}$.  A bad median requires at least
$m_\delta/2$ failures, so
\[
 \mathbb P\{\text{bad median}\}
 \leq\left(\frac5{4\sqrt2}\right)^{10s_\delta+1}
 \leq2^{-s_\delta}\leq\delta.
\]
The middle inequality follows from
$(5/(4\sqrt2))^{10}<1/2$.  Exact comparison of the rational outputs
and median selection have polynomial bit cost.

Finally, $N=10^{10}+n+\ell+\lceil\varepsilon^{-1}\rceil$ is bounded
by a fixed polynomial in the parameters of the theorem.  All construction,
evaluation, and counting bounds established above are fixed polynomials
in $N$ and the input encoding length.  Multiplying them by
$m_\delta=O(1+\log\delta^{-1})$ gives the required separate bounds on
rank-oracle calls and other bit operations on every execution, including
reading answers, writing queries and output, and generating random bits.
The only explicitly constructed labelled set has $2p$ elements, with
$p$ bounded by the short port ranges.  The large expansion matroids,
geometric subdivisions, and transport flows used in the proofs are
never constructed or queried.  This completes the algorithm and its
approximation and bit-complexity proofs.
\end{proof}
